\section{Results}

\paragraph{Motivating example.}
A deterministic \code{bias\_relu} template measured 1.029$\times$ above eager
in one fused8 screening run, then 0.976$\times$ at repeated median. Model
sampling cannot explain the reversal. We use it as an example, not an estimate
of how often such reversals occur. The new campaign assigns a distinct evidence
path to each research question (Table~\ref{tab:final-outcomes} and
Figure~\ref{fig:corrected-campaign}).

\begin{table}[t]
\centering
\small
\caption{Campaign summary. No KernelBench winner clears eager. Confirmation
removes one of three promotions in the parity stress test.}
\label{tab:final-outcomes}
\setlength{\tabcolsep}{3.5pt}
\renewcommand{\arraystretch}{1.10}
\begin{tabularx}{\linewidth}{@{}llrrrY@{}}
\toprule
Study & Evidence & Screen & Confirm & Dropped & Key estimate \\
\midrule
RQ1 holdout & 10 valid / 48 & 0 & 0 & -- & 1.014 ratio $[0.915,1.170]$ \\
RQ2 easy, $K=20$ & 4 tasks & 1.00 & 1.00 & 0 & $0.0000$ log optimism \\
RQ2 parity, $K=8$ & 4 tasks & 0.75 & 0.50 & 1/3 & $0.0072$ log optimism \\
RQ3 lifecycle & 24 rows & -- & -- & -- & 1.053 host $[0.998,1.114]$ \\
\bottomrule
\end{tabularx}
\end{table}

\paragraph{RQ1: validity precedes speed.}
We froze all 144 Gemini candidates before timing. Of the 141 records that reached
evaluation, 77 failed static policy and 9 failed correctness. Another 28 failed
during compilation or execution. The remaining 27 passed every gate. A compiler
OOM prevented evaluation of three candidates from one task. The valid candidates
covered 10 of 48 tasks, all matrix multiplications.

At screening, one selected winner beat \code{torch.compile} by 1.826$\times$
but remained below eager at 0.937$\times$. A distinct check across seven A4500
processes confirmed a 2.001$\times$ gain over the compiler. All 10 selected winners
remained below eager in screening and confirmation. The false promotion fraction
is therefore undefined, not zero. The median ratio from screening to confirmation
was 1.014 with task bootstrap interval $[0.915,1.170]$. This interval contains
parity and does not support directional selection optimism.

\paragraph{RQ2: candidate budget matters near parity.}
The easy regime acts as a negative control. All 80 deterministic candidates
confirmed above the margin. Screening and confirmation rates were 1.0 for every
$K\leq20$ on the original T4 and in a replication on the same A4500.

The calibrated stress test placed 32 frozen candidates near parity. At $K=8$,
screening promoted three of four task winners, while confirmation retained two.
The removed \code{bias\_gelu} winner moved from 1.0271$\times$ to
1.0001$\times$. The \code{bias\_relu} and \code{rmsnorm} winners remained above
the margin. Across candidate resamples, the screening and confirmation rates
diverged from 0.154 versus 0.124 at $K=1$ to 0.750 versus 0.500 at $K=8$. Median
log optimism at $K=8$ was 0.00725. Its task bootstrap interval,
$[-0.01239,0.02661]$, contains zero. The removed winner demonstrates the
mechanism but does not establish a population rate. Margin sensitivity appears
in Appendix~A.

\paragraph{RQ3: the timing boundary changes the result.}
The lifecycle control completed 24 process rows across eight tasks. Rebuilding
and transferring the reference inside the measured call raised synchronized
host latency by a median factor of 1.053. The enclosing CUDA event ratio was
1.0001. Their process row IQRs were $[0.998,1.114]$ and $[0.9919,1.0019]$.
Reference lifecycle work can therefore alter a host latency claim without
changing the enclosed device execution.

\begin{figure}[t]
  \centering
  \includegraphics[width=\linewidth]{corrected_campaign_results.pdf}
  \caption{Independent confirmation changes the supported claim. Panel A shows
  that all 10 valid KernelBench winners remain below the 1.02$\times$ eager
  margin. Panel B isolates one screening promotion that does not confirm. Panel
  C compares host and CUDA event ratios for the lifecycle control.}
  \label{fig:corrected-campaign}
\end{figure}
